Ambiguity check: the statement says ``intervals of consecutive integers''. If $0\in I_i$ for some $i$ then the product is $0$, which is a perfect power ($0=0^t$ for any $t\ge 2$). So the literal statement is false unless one restricts to (strictly) positive integers.

\medskip
\noindent\textbf{Corrected (standard) formulation.}
Fix an integer $r\ge 1$. An \emph{interval of consecutive positive integers} means a set
\[
I=[a,b]\cap\mathbb Z=\{a,a+1,\dots,b\}\subset\mathbb Z_{\ge 1}
\]
with $1\le a\le b$, and its \emph{length} is $|I|=b-a+1$.
Call a positive integer $P$ a \emph{perfect power} if $P=x^t$ for some integers $x\ge 2$ and $t\ge 2$.

The conjecture is:
\[
\forall r\ge 1\ \exists k\ge 1\ \forall\text{ pairwise disjoint intervals }I_1,\dots,I_r\subset\mathbb Z_{\ge1}
\ \Bigl(\min_i |I_i|\ge k\ \Rightarrow\ \prod_{i=1}^r\prod_{m\in I_i} m\ \text{is not a perfect power}\Bigr).
\]
